% TOP_PROBLEM: {"id": 30006857, "title": "Hartshorne Conjecture on complete intersections in codimension two", "category_id": 5, "category": {"id": 5, "name": "algebraic_geometry", "display_name": "Algebraic Geometry", "description": "Geometric objects defined by polynomial equations.", "slug": "algebraic-geometry", "order_index": 5, "created_at": "2026-07-31T15:26:25.671Z"}, "status": "open"}
% FIRST_POSTED: 2026-09-27
% !TeX program = lualatex
% ATTEMPT_STATUS: unresolved
\documentclass[11pt]{article}
\usepackage[margin=25mm]{geometry}
\usepackage{fontspec}
\setmainfont{Latin Modern Roman}
\usepackage{amsmath,amssymb,mathtools}
\usepackage{unicode-math}
\setmathfont{Latin Modern Math}
\usepackage{xurl,hyperref}
\hypersetup{colorlinks=true,urlcolor=blue}
\setcounter{secnumdepth}{0}
\setlength{\parindent}{0pt}
\setlength{\parskip}{5pt}
\setlength{\emergencystretch}{3em}
\begin{document}
\section*{\texorpdfstring{Hartshorne Conjecture on complete intersections in codimension two}{Hartshorne Conjecture on complete intersections in codimension two}}\label{30006857}
\subsection{Definitions and mathematical statement}
Let $X\subset{\mathbb{P}}^N_{{\mathbb{C}}}$ be a smooth projective subvariety of codimension two, with $N\ge7$. It is a scheme-theoretic complete intersection of two hypersurfaces if there are homogeneous polynomials $F,G$ of positive degrees such that the closed subscheme cut out by their ideal is exactly $X$:
\[
 X=V(F,G)\subset{\mathbb{P}}^N_{{\mathbb{C}}}.
\]
The codimension-two case of Hartshorne's conjecture asserts this for every such $X$. Scheme-theoretic equality records the defining ideal structure, not merely equality of complex point sets. No particular degrees for $F,G$ are prescribed, and the ambient-dimension threshold $N\ge7$ is part of the selected target. Smoothness refers to $X$, not necessarily to every hypersurface chosen to define it.
\par\smallskip\textit{Formulation and concept references:} \href{https://link.springer.com/article/10.1007/s12215-024-01032-4}{[S1]}.

\subsection{Short English statement}
Must every smooth codimension-two complex projective variety in projective space of dimension at least seven be cut out by two hypersurface equations?
\subsection{Research attempt: a smooth containing hypersurface and the rank-two splitting gap}
\paragraph{1. Scope and source audit.}
The conclusion is a scheme-theoretic intersection of two hypersurfaces over $\mathbb C$, with $X$ smooth of codimension two in $\mathbb P^N$ and $N\geq7$. Producing two equations with the same support would not suffice. Section 6 of [S1], checked on 27 September 2026, states precisely this codimension-two form and discusses subcanonicity via Barth--Larsen. The degree-dependent theorem of Erman, Sam, and Snowden [E1] applies when the ambient dimension is sufficiently large compared with the degree; it is not the uniform threshold $N\geq7$ for arbitrary degree. The arguments below use the standard Lefschetz theorem for the Picard group of a smooth hypersurface, Barth--Larsen, the Hartshorne--Serre construction, and Horrocks's splitting criterion as named external inputs. The reductions from these inputs are proved explicitly.

\paragraph{2. A complete subcase: containment in a smooth hypersurface.}
Suppose $X\subset H\subset\mathbb P^N$, where $H$ is a smooth hypersurface of degree $a>0$ and $X$ has codimension one in $H$. Since both are smooth, $X$ is an effective Cartier divisor on $H$. By the Picard form of Lefschetz, valid here because $\dim H=N-1\geq6$,
\[
 \operatorname{Pic}(H)=\mathbb Z[\mathcal O_H(1)].
\]
Write $\mathcal O_H(X)=\mathcal O_H(b)$. The divisor is nonempty, so intersecting it with $N-2$ general hyperplanes gives $\deg X=ab>0$, and hence $b>0$.

Let $F$ define $H$. The canonical section of $\mathcal O_H(b)$ vanishing on $X$ lifts to a degree-$b$ homogeneous polynomial $G$ on $\mathbb P^N$, because the sequence
\[
 0\longrightarrow\mathcal O_{\mathbb P^N}(b-a)
 \xrightarrow{\cdot F}\mathcal O_{\mathbb P^N}(b)
 \longrightarrow\mathcal O_H(b)\longrightarrow0
\]
has $H^1(\mathbb P^N,\mathcal O(b-a))=0$. This is the standard intermediate-cohomology vanishing for line bundles on projective space, and it holds for every integer $b-a$ when $N\geq2$.

On $H$, the section $G|_H$ generates the ideal of the Cartier divisor $X$. Therefore the ideal sheaf of $X$ in $\mathbb P^N$ is generated by $F$ and $G$, not merely radical-equivalent to their ideal. Thus $X=V(F,G)$ scheme-theoretically. This proves the desired conclusion whenever a smooth containing hypersurface exists.

The gap is not the existence of some containing hypersurface: every projective subvariety has many equations. It is the existence of one that is smooth, particularly along $X$. Bertini applied to a linear system with base locus $X$ gives smoothness away from that base locus; it does not by itself give smoothness along $X$.

\paragraph{3. The missing condition in a Bertini shortcut.}
For a homogeneous polynomial $F\in H^0(\mathcal I_X(a))$, its first derivative normal to $X$ is a section
\[
 dF|_X\in H^0\bigl(X,N^*_{X/\mathbb P^N}(a)\bigr).
\]
At a point of $X$, tangent derivatives vanish because $F$ restricts to zero on $X$. Thus the hypersurface $V(F)$ is smooth at that point exactly when this conormal section is nonzero there. To obtain smoothness all along $X$, one needs a nowhere-zero section of this rank-two bundle. Global generation, when available after a high twist, does not guarantee a nowhere-zero section: a general section of a rank-two bundle is expected to have a codimension-two zero locus. That expected locus can have positive dimension because $\dim X=N-2\geq5$.

This also explains why taking a very high-degree equation is not an automatic cure. Positivity can make the bundle generated, while its top Chern class obstructs a nowhere-zero section. A proof that one suitable normal derivative never vanishes would already contain substantial geometry absent from the simple Bertini argument.

\paragraph{4. The rank-two vector bundle reduction.}
For the present range, Barth--Larsen gives $\operatorname{Pic}(X)=\mathbb Z[\mathcal O_X(1)]$. In particular
\[
 \omega_X\cong\mathcal O_X(e)
\]
for an integer $e$. The Hartshorne--Serre correspondence, using the codimension-two local complete intersection and the projective-space cohomology vanishings, supplies a rank-two vector bundle $E$ on $\mathbb P^N$ and a section with zero scheme $X$, fitting into
\[
 0\longrightarrow\mathcal O_{\mathbb P^N}
 \longrightarrow E\longrightarrow\mathcal I_X(c)\longrightarrow0,
 \qquad c=e+N+1,\qquad\det E=\mathcal O(c).
\]
This existence statement is a standard external theorem, not an elementary consequence of the degree of $X$.

If $E\cong\mathcal O(a)\oplus\mathcal O(b)$, its section is a pair of homogeneous polynomials of degrees $a,b$. Because its zero scheme has pure codimension two and is nonempty, neither component can vanish identically, and neither degree can be zero or negative. A negative-degree component is necessarily zero; a nonzero degree-zero component has no zero; an identically zero degree-zero component leaves at most one equation. All these alternatives contradict the given zero scheme. Hence $a,b>0$, and the section gives the desired scheme-theoretic complete intersection.

Conversely, a complete intersection of degrees $a,b$ is the zero scheme of the tautological section of $\mathcal O(a)\oplus\mathcal O(b)$. Its Koszul resolution gives
\[
 0\longrightarrow\mathcal O(-a-b)
 \longrightarrow\mathcal O(-a)\oplus\mathcal O(-b)
 \longrightarrow\mathcal I_X\longrightarrow0.
\]
The vector-bundle reduction therefore identifies the difficult step as splitting the relevant rank-two bundle, rather than simply constructing a bundle or a section.

\paragraph{5. A conditional proof under an explicit cohomological hypothesis.}
Assume, in addition, that $X$ is arithmetically Cohen--Macaulay in the sheaf-cohomological sense
\[
 H^j(\mathbb P^N,\mathcal I_X(t))=0
 \quad\text{for all }t\in\mathbb Z\text{ and }1\leq j\leq N-2.
\]
Twist the Serre sequence by $\mathcal O(t)$. Since line bundles on $\mathbb P^N$ have no intermediate cohomology,
\[
 H^j(\mathbb P^N,E(t))\cong
 H^j(\mathbb P^N,\mathcal I_X(c+t))=0
 \qquad(1\leq j\leq N-2).
\]
The remaining intermediate degree follows by Serre duality and the rank-two identity $E^*\cong E(-c)$:
\[
 H^{N-1}(E(t))^*\cong H^1(E^*(-t-N-1))
 \cong H^1(E(-c-t-N-1))=0.
\]
Horrocks's criterion now says that $E$ splits as a sum of line bundles. Section 4 then proves that $X$ is a complete intersection. This is a complete deduction from the explicit ACM hypothesis and the named standard theorems.

The hypothesis includes $H^1(\mathcal I_X(t))=0$ for every twist. It must not be shortened to only the higher intermediate cohomology of $\mathcal O_X$. Smoothness implies local Cohen--Macaulayness, whereas ACM is a global condition on the homogeneous coordinate ring and its deficiency modules. Confusing these two notions would make the argument appear to solve the conjecture by deleting its essential input.

\paragraph{6. Numerical checks are necessary but not a splitting theorem.}
For a complete intersection,
\[
 \deg X=ab,\qquad e=a+b-N-1,\qquad
 c_1(E)=a+b,\quad c_2(E)=ab.
\]
Thus a proposed splitting must satisfy
\[
 c_1(E)^2-4c_2(E)=(a-b)^2,
\]
a nonnegative square integer, with the parity needed to recover integral $a,b$. Failure of this test rules out a proposed pair of degrees. Passing it does not prove that the bundle splits or that two equations generate the ideal. Chern classes contain less information than an isomorphism class of vector bundles.

One must also distinguish a degree-dependent theorem from the conjectured fixed codimension threshold. If a result requires $N\geq B(\deg X)$, it supplies a genuine subcase, but an arbitrary degree cannot be inserted into a fixed value such as $N=7$. The degree may grow without bound in the question.

\paragraph{7. Why raising ambient dimension by cones does not give a counterexample.}
The twisted cubic in $\mathbb P^3$ is a smooth codimension-two curve of degree three and is not a complete intersection. Indeed, it lies in no hyperplane, so if it were an intersection of two hypersurfaces their degrees would both be at least two, forcing degree at least four. Its ideal is generated by quadratic minors and has no linear equation.

Taking a projective cone embeds related codimension-two schemes in larger projective spaces, but the vertex is singular. In affine coordinates centered at a vertex point, the defining equations have no linear terms, so the Zariski tangent space is the full ambient tangent space. The cone itself has codimension two. Its tangent dimension therefore exceeds its local dimension, proving singularity. Iterating the cone operation preserves this defect. Hence these low-dimensional examples do not violate the smooth high-dimensional statement.

\paragraph{8. Diagnostics and outcome.}
The exact diagnostic enumerates positive degree pairs in a finite range and verifies the adjunction/discriminant identities, as well as the elementary degree obstruction for the twisted cubic. It does not compute an arbitrary bundle's splitting type or test ACM for all twists. Those universal conditions are addressed only in the conditional proofs.

\textbf{UNRESOLVED.} Smooth hypersurface containment and the explicit ACM condition each imply the required scheme-theoretic conclusion. The first unrestricted gap is to establish such a condition, or directly split the associated rank-two bundle, for every smooth codimension-two $X\subset\mathbb P^N$ with $N\geq7$. Neither local smoothness, Bertini away from the base locus, nor compatible Chern numbers accomplishes that step.

\subsection{Sources}
\begin{itemize}
\item[S1]Homemade algebraic geometry, Section 6, codimension-two restatement.
\url{https://link.springer.com/article/10.1007/s12215-024-01032-4}.
\textit{Formulation source.}
\item[E1]D. Erman, S. V. Sam and A. Snowden, Strength and Hartshorne's Conjecture in high degree, arXiv:1804.09730.
\url{https://arxiv.org/abs/1804.09730}.
\textit{Primary abstract checked 27 September 2026: a dimension bound depending on degree, not the fixed threshold of the supplied target. The full paper was not independently audited.}
\end{itemize}
\end{document}
